% Template SCALE LAC conference

\documentclass[a4paper]{article}
\usepackage[utf8]{inputenc}
\usepackage{amsmath}
\usepackage{amssymb,amsfonts,textcomp}
\usepackage[T1]{fontenc}
\usepackage[english,spanish]{babel}
\usepackage{color}
\usepackage{array}
\usepackage{hhline}
\usepackage{hyperref}
\usepackage[pdftex]{graphicx}
\usepackage{lipsum}  

% Outline numbering
\setcounter{secnumdepth}{2}

% Page layout (geometry)
\setlength\voffset{-1in}
\setlength\hoffset{-1in}
\setlength\topmargin{1.0236in}
\setlength\oddsidemargin{0.9839in}
\setlength\textheight{9.0164995in}
\setlength\textwidth{6.3002996in}
\setlength\footskip{0.4327in}
\setlength\headheight{0.2362in}
\setlength\headsep{0.1965in}


% Pages styles
\makeatletter
\newcommand\ps@Standard{
\renewcommand\@oddhead{}
\renewcommand\@evenhead{}
\renewcommand\@evenfoot{\thepage}
\renewcommand\@oddfoot{\thepage}
\renewcommand\thepage{\arabic{page}}
}

% HEADER and FOOTER
\newcommand\ps@FirstPage{
\renewcommand\@oddhead{\itshape\small MIT SCALE Conference for Latin America and the Caribbean ${\bullet}$  Cambridge, MA, U.S.  ${\bullet}$  March 20th -23th, 2022}
\renewcommand\@evenhead{\@oddhead}
\renewcommand\@oddfoot{\itshape Paper title here\hfill \thepage}
\renewcommand\@evenfoot{\@oddfoot}
\renewcommand\thepage{\arabic{page}}
}



\begin{document}

\clearpage\setcounter{page}{1}%\pagestyle{Standard}
\thispagestyle{FirstPage} % header and footer first page

%%%%%%%%%%%%%%%%%%% Title %%%%%%%%%%%%%%%%%%%%

\begin{center}
	\begin{LARGE}
		\textbf{Short Paper's Title}
	\end{LARGE}
\end{center}

\bigskip

%%%%%%%%%%%% List of authors %%%%%%%%%%%%%%%

{\centering
	Author 1\textsuperscript{1}, Author 2\textsuperscript{2}, Author 3\textsuperscript{2} and Author 4\textsuperscript{3
	}\textit{(e.g., Sampaio, N.)}\par}
\bigskip
{\centering
	1: Affiliation Department, University, Country 	 \par}
{\centering
	2: Affiliation Department, University, Country   \par}
{\centering
	3: Affiliation Department, University, Country   \par}

\bigskip

%%%%%%%%%%%% Abstract %%%%%%%%%%%%
{\centering\bfseries\color{black}Abstract 
	\par}
Abstract should contain between 150 and 200 words long, summarizing the main idea of the problem or research question, the methodology was used, the key contributions and the main findings.
\textcolor{red}{\lipsum[1-1]}

\bigskip

%%%%%%%%%%%% Keywords %%%%%%%%%%%%
\textit{\textbf{Keywords}: List here at most five keywords of your research }

\bigskip

{\centering\bfseries\color{red} (Following content must be between 2200 and 3100 words plus tables, charts and figures.) \par}% Remove this line at the end

%%%%%%%%%%%% Body %%%%%%%%%%%%

\section[Introduction]{ Introduction}
It includes statistics/facts, justification and motivation about the study. It must contain goal, research question and
scope of the research. You should provide an overview of the methodology and the most important findings.
\textcolor{red}{\lipsum[1-1]}

\section[Literature review]{Literature review}

It describes the state-of-the-art body and analyzes the most representative studies related to the research project. It
must include a summary of methodological and practical gaps identified by authors that clearly positions the main
contributions of the research.
\textcolor{red}{\lipsum[1-1]}


\section[Methods and procedures]{ Methods and procedures}
This section provides a description of the theoretical framework you applied to do your research, the main methods and
techniques used to collect/process data, perform analysis, etc. Document mathematical formulations and solution
algorithms. Please number formulas to ease explanation elsewhere in the manuscript by citing the numbers of equations,
inequalities, etc.
\textcolor{red}{\lipsum[1-1]}


\section[Experimental/numerical setting]{ Experimental/numerical setting}
The experimental setting, the values of parameters or the data collection to feed a practical/empirical model (connected
to a case study) are described in detail in this section. The description of the numerical design or data set must be
linked with the methodological framework and the model introduced in the previous section. If you work with a company,
organization make sure to describe their problem and case study in a concise, clear way.
\textcolor{red}{\lipsum[1-1]}


\section[Results and Discussion]{\rmfamily Results and Discussion}

Please explain in detail your results and discuss main findings, recommendations, managerial insights and most important
methodological results. Make use of Figures and tables to ease explanation wherever is needed (i.e. Figure~\ref{fig}).
\textcolor{red}{\lipsum[1-1]}


\begin{figure}
	\centering\includegraphics[width=2in,height=2.3819in]{test.png}
	\caption{Image}
	\label{fig}
\end{figure}

\section[Conclusions and future research]{Conclusions and future research}

Please add your main conclusions here and future research venues
\textcolor{red}{\lipsum[1-1]}


\section[Acknowledgements (optional)]{Acknowledgements (optional)}

It includes grants, acknowledgements to people, institutions, etc.
\textcolor{red}{\lipsum[1-1]}


\section[References]{References}
Use Reference Style `Science' and follow APA version 7.0 rules on your citations and references.


\section[Appendices (optional)]{ Appendices (optional)}
Add all supplementary material such as extended formulations, theorems, proofs, raw data, experiments, nomenclature in
this section. 


\textcolor{blue}{The paper must be 5 to 7 pages long including tables, charts and figures. The file size must be no more than 10MB.}

\end{document}
